\documentclass[12pt]{article}
\usepackage[utf8]{inputenc}
\usepackage[T1]{fontenc}
\usepackage{hyperref}
\title{Central Bikol Language Resources for K-12 Curriculum Development}
\author{Compiled by Claude Sonnet 5.5 (claude-sonnet-5-5) for Andr\'as Kornai}
\date{}
\begin{document}
\maketitle

\section{Introduction}
Central Bikol (ISO 639-3 \texttt{bcl}) is spoken mainly in Naga City, Tabaco, Legazpi and Sorsogon City in the Philippines. Bikol Naga was assigned as the L1 for mother-tongue-based multilingual education (MTB-MLE), although the dialects of the six Bikol provinces differ noticeably. The Canaman dialect is the standard form in literature, Catholic rites and mass media.

\section{Primary Educational Resources}
\begin{itemize}
\item DepEd LRMDS, ``MTB-MLE Grade 1: Bikol Week 30 Teacher's Guide'' (2013): \url{https://lrmds.deped.gov.ph/detail/3961}
\end{itemize}
The LRMDS portal is the most direct place to look for further grade-level Bikol materials; I did not find other specific titles in search results.

\section{Linguistic and Reference Materials}
\begin{itemize}
\item Mintz, Bikol Grammar Notes (University of Hawai`i Press, open access): \url{https://manifold.uhpress.hawaii.edu/projects/bikol-grammar-notes}
\item Wikipedia, ``Central Bikol'': \url{https://en.wikipedia.org/wiki/Central_Bikol}
\item Wikipedia, ``Bikol languages'': \url{https://en.wikipedia.org/wiki/Bikol_languages}
\item Bicol University journal article (Lorenzana): \url{https://systems.bicol-u.edu.ph/journal/assets/journal_pdf/4%20Lorenzana_49-61.pdf}
\item Bicol University journal article (de la Torre): \url{https://systems.bicol-u.edu.ph/journal/assets/journal_pdf/1_de%20la%20Torre_1-12.pdf}
\end{itemize}
The Mintz materials include a fairly comprehensive dictionary, aimed at Peace Corps learners rather than schools.

\section{Religious and Community Resources}
The Global Recordings Network page lists a Central Bikol New Testament audio recording from Faith Comes By Hearing: \url{https://globalrecordings.net/en/language/bcl}. A full printed Bible translation was not clearly documented in my searches.

\section{Digital and Audio Resources}
See the Global Recordings Network page above. No dedicated digital dictionary or corpus was found in search results.

\section{Curriculum Development Context}
Central Bikol is used as a language of instruction under MTB-MLE. A 2023 Bicol University study notes that the main problems are limited awareness of and access to existing written and recorded materials, and weak support for producing literature. Dialect variation across provinces is a standing challenge for standardising materials.
\end{document}